\documentclass[10pt,a4paper]{amsart}
\usepackage[margin=19mm]{geometry}
\usepackage{amsmath,amssymb,mathtools}
\usepackage[scr=rsfs,cal=euler]{mathalfa}
\usepackage{xfrac}
\usepackage[unicode,hidelinks,bookmarks=false]{hyperref}
\newcommand{\ie}{i.e.\ }
\newcommand{\cf}{cf.\ }
\newcommand{\resp}{respectively\ }
\newcommand{\Subsection}{\subsection}
\providecommand{\nobreakdash}{\nobreak\hbox{-}}
\setlength{\emergencystretch}{2em}
\multlinegap0pt
\usepackage{amssymb}
\usepackage{enumerate}
\usepackage[shortlabels]{enumitem}
\usepackage{cancel}
\usepackage{tikz}
\usepackage{xfrac}
\usepackage{algorithm}\usepackage{algpseudocode}
\usepackage{mathtools}
\newtheorem{lemma}{Lemma}
\newtheorem{theorem}{Theorem}
\usepackage{cleveref}
\crefname{lemma}{Lemma}{Lemmas}
\crefname{theorem}{Theorem}{Theorems}
\newcommand{\Ical}{\mathcal{I}}
\newcommand{\R}{\mathbb{R}}
\newcommand{\inner}[2]{\left\langle {#1}, {#2} \right\rangle}
\title{Inexactly Smooth Performance Estimation\\ and New Optimized Gradient Methods}
\author{Aaron Zoll, Benjamin Grimmer}
\date{Source: arXiv:2606.01505v2 (CC BY 4.0)}
\begin{document}
\maketitle

\noindent\emph{Source and licence.} Inexactly Smooth Performance Estimation\\ and New Optimized Gradient Methods. Version arXiv:2606.01505v2 (CC BY 4.0), distributed by arXiv under the Creative Commons Attribution 4.0 International licence, http://creativecommons.org/licenses/by/4.0/, as recorded in the arXiv licence metadata for this exact version and shown on the abstract page https://arxiv.org/abs/2606.01505v2. Full licence notice: \texttt{licenses/arxiv-2606.01505-CC-BY-4.0.txt}.\par\smallskip

\noindent\textit{Editorial scope.} Counted unit: the article's theorem thm:interpolation-general (source0.tex lines 491--497). The article's complete original proof of it is delivered with it (source0.tex lines 588--607, one \texttt{proof} environment of 20 lines, which the article places away from the statement). No mathematics is written, repaired or re-derived: the statement and the proof are the article's own lines, in the article's own notation, and no retained line is substituted or re-flowed.  Retained uncounted as the source-local context the proof needs: lines 366--373; lines 388--395; lines 399--401.  Nothing was omitted from any retained span, so the package carries no omission record.  No retained line contains a citation, so the wrapper carries no bibliography and no bibliography entry.  No retained line contains a figure environment, an image-inclusion command or drawing code, so no image is delivered and the package declares no asset dependency.  Editorial formatting only, in addition: an AMS \texttt{amsart} preamble that restates the article's own declarations and macros used by the retained lines, namely the environments \texttt{lemma}, \texttt{theorem} on separate counters, together with the standard packages those lines need.  The retained environments print in this document as Lemma 1, Theorem 1 in order of appearance, whereas the article prints the same items with the numbering of its own full document.  The complete native source of the article as distributed in the arXiv e-print for this exact version is bundled unchanged as \texttt{sources/201834/source0.tex}.
\begin{lemma}\label{lemma:L(delta) conjugate equivalence}
    Let $f$ be closed, convex, and proper. The following are equivalent:
    \begin{enumerate}        
        \item $f(y) \geq f(x) + \inner{g_x}{y-x}+\frac{1}{2L(\delta)}\|g_y-g_x\|^2 - \delta, \quad \forall \delta \geq 0, \ \forall x,y \text{ with } g_x \in \partial f(x),\ g_y\in \partial f(y)$,
        \item $f$ is $L(\cdot)$-inexactly smooth,
        \item $f^*$ is $1/L(\cdot)$-inexactly strongly convex. %$f^*(g_y) \geq f^*(g_x)+ \inner{x}{g_y-g_x}+ \frac{1}{2L(\delta)}\|g_x-g_y\|^2 -\delta, \quad \forall g_x, g_y \text{ with }x \in \partial f^*(g_x), y \in \partial f^*(g_y)$.
    \end{enumerate}
\end{lemma}

\begin{lemma}\label{lemma:L(delta) sums and max calculus}
    Consider any closed, convex, and proper $L_i(\cdot)$-inexactly smooth functions $f_i$ for $i=1,\dots,m$. Then, \begin{enumerate}
        \item $\sum_{i=1}^m f_i$ is $\left(\inf_{\substack{\sum_{i=1}^m \delta_i = \delta\\ \delta_i \geq 0}} \sum_{i=1}^m L_i(\delta_i)\right)$-inexactly smooth,
        \item $(\max_i f_i^*)^*$ is $\max_i L_i(\delta)$-inexactly smooth,
        \item $\left(\sum_{i=1}^m f_i^*\right)^*$ is $\left(\sup_{\substack{\sum_{i=1}^m \delta_i = \delta\\ \delta_i \geq 0}} \sum_{i=1}^m \frac{1}{L_i(\delta_i)}\right)^{-1}$-inexactly smooth.
        %\item $\displaystyle\max_i f_i$ is $\left(\displaystyle\inf_{\substack{\delta_1 + \delta_2 = \delta\\ \delta_1, \delta_2 \geq 0}} \left\{\max_i L_i(\delta_1)+\frac{M^2}{2\delta_2}\right\}\right)$-inexactly smooth where $$M = \sup_x \{\|g_x^{(i)} -g_x^{(j)}\| : g_x^{(i)} \in \partial f_i(x), \  g_x^{(j)} \in \partial f_j(x)\}. $$
    \end{enumerate}
\end{lemma}

\begin{lemma}\label{composing with abs to norm}
    Consider any convex function $h : \R \to \R $ such that $h(|t|)$ is $L(\cdot)$-inexactly smooth. Then $f(x) = h(\|x\|)$ is $L(\cdot)$-inexactly smooth.  
\end{lemma}

\begin{theorem} \label{thm:interpolation-general}
   Let $L(\cdot)$ satisfy our assumptions. If a set of observations $\mathcal{H} =\{(x_i,f_i,g_i)\}_{i\in\Ical}$ is $L(\cdot)$-interpolable, then for each $i,j\in\Ical$,
\begin{equation}\label{eq:interpolation condition}
        f_i - f_j - \langle g_j, x_i-x_j\rangle - \frac{1}{2L(\delta)}\|g_i-g_j\|^2 + \delta \geq 0, \quad \forall \delta \geq 0 . 
    \end{equation} 
    Conversely, if the above condition holds then  $\mathcal{H}$ is ${c_L L(\cdot/c_L)}$-interpolable.
\end{theorem}

\begin{proof}[{\bf Proof of \cref{thm:interpolation-general}}]
Suppose $f(x)$ is an $L(\cdot)$-inexactly smooth function that interpolates the observations $\mathcal{H}$. By the cocoercive-type condition 1.~from \cref{lemma:L(delta) conjugate equivalence} with $x = x_j$ and $y = x_i$, the claimed conditions hold for each $i, j \in \Ical$. 

For the reverse direction, define the function $r = (\max_{i\in \Ical} r_i^*)^*$ where
\begin{align*}
     r_i(x) &= f_i + \inner{g_i}{x-x_i}+\theta_L(\|x-x_i\|).
 \end{align*} 
 Calculating the conjugates of each $r_i$, one has $r_i^*(g) = -f_i + \inner{x_i}{g}+\theta_L^*(\|g-g_i\|) $.
 
 This claimed interpolation $r$ is convex as all conjugates are convex. Our calculus rules allow us to verify its inexact smoothness. Each individual $r_i$ is $c_L L(\cdot/c_L)$-inexactly smooth by the sum rule in Lemma~\ref{lemma:L(delta) sums and max calculus} since the linear term $f_i + \inner{g_i}{x-x_i}$ does not affect the smoothness and $\theta_L(\|x-x_i\|)$ is $c_L L(\cdot/c_L)$-inexactly smooth by Lemma~\ref{composing with abs to norm}. Then the $c_L L(\cdot/c_L)$-inexact smoothness of $r$ follows from the conjugate maximum formula in Lemma~\ref{lemma:L(delta) sums and max calculus}.

 All that remains is to verify that $r$ interpolates the given first-order information. Note that $r(x_j) = f_j$ and $g_j\in\partial r(x_j)$ if and only if $\max_{i\in \Ical} r_i^*(g_j) = -f_j+\langle x_j,g_j\rangle$ and $x_j \in \partial (\max_{i\in \Ical} r_i^*)(g_j)$. The key step in establishing this is showing that the interpolation conditions guarantee that $r_j^*$ attains this maximum at $g_j$. To see this, observe that for any $i\in\Ical$,
 \begin{align*}
     r_j^*(g_j) &= -f_j + \inner{x_j}{g_j}
     \geq -f_i + \inner{x_i}{g_j}+\theta_L^*(\|g_i-g_j\|)
     =r_i^*(g_j),
 \end{align*}
 where the inequality is our interpolation condition between $i$ and $j$ (taking the supremum over $\delta \geq 0$).
 From this, the fact that $r$ interpolates follows since $r^*(g_j)=-f_j+\langle x_j,g_j\rangle$ and $x_j \in \partial r_j^*(g_j)\subseteq \partial (\max_{i\in \Ical} r_i^*)(g_j)$.
\end{proof}

\end{document}
